\paragraph{Soporte celular del transporte central.}
La recurrencia anterior determina también la geometría finita de la
trayectoria. Consideramos las posiciones inmediatamente anteriores a cada
uno de los \(108\) pasos; las coordenadas pertenecen a
\(D_9=\{1,\ldots,9\}\), con transporte cíclico módulo nueve.

\begin{proposition}[Tres diagonales cíclicas de la sección central]
\label{prop:el-soporte-27}
La proyección celular de la trayectoria iniciada en
\((r,c,h,v)=(5,5,1,1)\) tiene soporte exacto
\begin{equation}
 \begin{aligned}
 \operatorname{Supp}_{\mathrm{cell}}(\Gamma_e)&=\{(r,c)\in D_9^2:
                 c-r\equiv0,1,-1\pmod9\},
 \\
 |\operatorname{Supp}_{\mathrm{cell}}(\Gamma_e)|&=27.
 \end{aligned}
 \label{eq:el-soporte-27}
\end{equation}
La orientación inicial conjugada \((h,v)=(-1,-1)\), con el mismo
calendario de inversiones, produce el mismo soporte.
\end{proposition}
\begin{proof}
Dentro de un bloque de \(27\) pasos, escribimos \(h=v=s\), donde
\(s\in\{-1,1\}\). Si un bloque de tres pasos comienza en \((a,a)\),
la lectura aditiva precede al desplazamiento hasta \((a,a+s)\);
la lectura multiplicativa precede al desplazamiento hasta
\((a+s,a+s)\); el umbral conserva esta última posición. Los tres
estados previos pertenecen, por tanto, a las diagonales
\(c-r=0\) y \(c-r=s\) del toro posicional. Los nueve bloques
consecutivos recorren las nueve posibilidades de \(a\), pues la
traslación por \(s\) tiene orden nueve. La inversión siguiente
recorre la diagonal adyacente de signo contrario. Las diagonales
de diferencias \(0,1,-1\) son disjuntas, cada una contiene nueve
celdas y todas quedan visitadas. Intercambiar el signo inicial
intercambia el orden de visita de las dos diagonales adyacentes.
\end{proof}

El punto fijo \((5,5)\) caracteriza la sección de origen; el conjunto
\(\operatorname{Supp}_{\mathrm{cell}}(\Gamma_e)\) caracteriza su trayectoria celular. El resultado
determina una morfología discreta sobre el dominio APP. La realización
espacial y el pegado de orientación reciben además los datos de
transporte que les corresponden. La permanencia de la sección y la
variación de su posición observable son así propiedades compatibles.

\paragraph{Dos historias centrales en la coordenatización integral.}
La coordenatización \(\mathcal L_{12}\) del
apartado~\ref{subsec:k-registro-integral} conserva los eventos y su
orden. La aplicamos a las dos orientaciones centrales. Para cada
paso \(t\), sea \(d_t=\rho_9(u_t)\), donde \(u_t\) es el entero
aditivo, multiplicativo o de umbral leído por la recurrencia.
En la ventana \(E_{i+1}=\{9i+1,\ldots,9i+9\}\), definimos
\[
 n_i=\#\{t\in E_{i+1}:d_t\in\{2,4,6,8\}\},\qquad
 e_i=\Sigma_{12}(i+1)n_i,\qquad0\le i<12.
\]
El signo \(\Sigma_{12}\) es la orientación sectorial de la coordenatización.
En la trayectoria \({++}\) coincide con el signo horizontal
transportado; en la trayectoria \({--}\) coincide con su opuesto.
Equivalentemente, en ambas es
\(\Sigma_{12}(i+1)=h^{\mathrm{pre}}_{9i+1}/h_0\), donde
\(h^{\mathrm{pre}}_t\) es la orientación anterior a la lectura del paso.
Esta normalización conserva por separado la orientación inicial y
la orientación relativa. Los símbolos \(e_i\) designan eventos
enteros, mientras que \(e\) sin índice designa el número de Euler.

La evaluación por ventanas de las dos recurrencias es
\[
\begin{array}{c|rrrrrrrrrrrr}
i&0&1&2&3&4&5&6&7&8&9&10&11\\\hline
n_i^{++}&2&2&5&4&3&2&2&2&5&4&3&2\\
n_i^{--}&4&3&2&2&2&5&4&3&2&2&2&5
\end{array}
\]
y determina las historias
\begin{equation}
 e_{++}=(2,2,5,-4,-3,-2)^2,\qquad
 e_{--}=(4,3,2,-2,-2,-5)^2.
 \label{eq:el-historias-centrales}
\end{equation}
La potencia indica concatenación. Ambas historias tienen suma cero,
la misma sucesión de signos sectoriales y el mismo histograma
absoluto: seis doses, dos treses, dos cuatros y dos cincos.

\begin{proposition}[Separación reversible de las dos historias]
\label{prop:el-memorias-distintas}
En las coordenadas \((s_C,M,a)=\mathcal L_{12}(e)\), ambas historias
satisfacen \(s_C=0\) y \(a=0\), mientras que
\begin{equation}
 M_{++}=(8,8,14,-4,-2),\qquad
 M_{--}=(18,16,14,6,6).
 \label{eq:el-memorias-distintas}
\end{equation}
La coordenatización reconstruye cada historia de manera única.
\end{proposition}
\begin{proof}
Las dos sextuplas se repiten, de modo que
\(p_i=e_i+e_{i+6}=2e_i\) y \(a_i=e_i-e_{i+6}=0\).
Restar \(p_5\) a los primeros cinco \(p_i\) produce los márgenes
de \eqref{eq:el-memorias-distintas}. Sus sumas son \(24\) y
\(60\), respectivamente; la fórmula inversa del
lema~\ref{lem:k-carta-integral} da \(p_5=-4\) y \(p_5=-10\).
Después, \(p_i=p_5+M_i\) y
\(e_i=e_{i+6}=p_i/2\) recuperan exactamente las dos sextuplas.
Se satisfacen la divisibilidad por seis y las seis congruencias de
paridad que caracterizan la imagen integral del lector.
\end{proof}

La igualdad de censos y signos conserva una proyección común; los
márgenes conservan la diferencia de historia. La recurrencia opuesta
emite \(\gamma_{--}=(b^3a^3)^2\): tras \(27\) pasos la ruta
\({++}\) alcanza de nuevo el centro con la orientación inicial de
\({--}\), y el calendario de tres pasos conserva su fase. Así,
\(\gamma_{--,t}=\gamma_{e,t+27}\), con índices cíclicos.
Los recuentos \(D_k\) son invariantes bajo ese desplazamiento; el
lector de ventanas recibe la misma palabra sectorial en las dos
historias. Ambos lectores de retorno coinciden, por tanto, para
estos registros que la coordenatización integral distingue. La evaluación másica y
la reconstrucción de eventos poseen alcances informacionales
precisos dentro del mismo transporte.
